\documentclass[11pt,letterpaper]{article}
\usepackage{hanzo-defense}

\doctitle{Cloud Economics of a Sovereign OSS Stack:\\A Regulated Capital-Markets\\Migration Case Study}
\docid{HAI-2026-012}
\docversion{Version 1.0}
\docdate{2026}
\docauthors{Zach Kelling --- Hanzo Team}
\docorgs{Hanzo Industries\\research@hanzo.ai}

\begin{document}
\makecover

\begin{abstract}
We document the infrastructure economics of migrating a regulated capital-markets
platform from a sprawling multi-project Google Cloud (GCP) footprint to a
consolidated, self-hostable open-source stack built with Hanzo's agentic
engineering system. The client---a FINRA/SEC-regulated securities platform,
anonymized here---was spending, by \emph{actual billing} once a startup-credit
program expired, approximately \$35{,}158/month ($\approx\$421{,}900$/year) across
eleven billing accounts, the great majority consumed by legacy GKE clusters and GPU
virtual machines left over from an earlier architecture. Consolidation onto the new
stack and owned hardware reduced spend to approximately \$3{,}500/month
($\approx\$42{,}000$/year)---a $\sim$90\% reduction, on the order of \$380{,}000/year
saved---and removed the dependency on a single cloud vendor. The same agentic stack
also \emph{built} the platform: a full SEC-registered ATS with broker-dealer and
transfer-agent functions, consolidated from 200+ microservices and
3{,}071{,}930 active lines of a prior implementation down to 1{,}056{,}566 active
lines ($\approx$66\% less code; $\sim$2.0\,M lines of dead code deleted) running as
3 compiled Go binaries plus a small set of OSS platform services, carried through
FINRA/SEC approval and SOC\,2. We decompose where the savings come from, report the
figures with their provenance, and argue the same mechanism applies directly to
any regulated enterprise or government modernization---defense among them---
carrying a metered-API-plus-hyperscaler-sprawl cost structure.
\end{abstract}

\paragraph{Executive Summary.}
A regulated financial client appeared to be spending only a modest amount on Google
Cloud because a startup-credit program was absorbing the bill. When those credits
expired, the \emph{actual} cost surfaced: about \$35{,}158 a month---roughly \$926
a day of newly cash-billed spend---almost all of it on old GKE clusters and GPU
machines the new architecture had made unnecessary. Migrating to a consolidated,
self-hostable open-source stack on owned hardware cut that bill to about \$3{,}500 a
month ($\approx$90\%, on the order of \$380{,}000 a year), and removed the
dependency on a single cloud vendor; ongoing consolidation targets a further floor
near \$1{,}200/month (nearly 97 percent off the original bill) as the last legacy
projects are retired. Just as
important is \emph{how little code} now runs that platform: the agentic engineering
stack collapsed 200+ legacy microservices into 3 compiled Go binaries plus a small
set of OSS platform services, cutting the active codebase from $\sim$3.07\,M to
$\sim$1.06\,M lines (deleting roughly 2.0\,M lines of dead code and eliminating
1.3\,GB of vendored dependencies)---and still cleared the bar of a heavily
regulated, audited domain (FINRA/SEC approval, SOC\,2). For any organization
carrying a similar cost structure---a healthcare provider, a financial firm, a
critical-infrastructure operator, or a government contractor---the lesson is
concrete: the same approach modernizes an existing stack to a sovereign,
self-hosted footprint that is far cheaper
to run, smaller to maintain, and faster to build.

\tableofcontents
\newpage

% ============================================================================
\section{Before: A Sprawling, Credit-Masked Footprint}
% ============================================================================

The starting state is typical of an organization that grew fast on Google Cloud
under promotional credits. Headline spend looked small because a startup-credit
program was absorbing it; an audit of the \emph{actual} GCP billing console---after
the credits expired---put true spend at approximately \textbf{\$35{,}158/month}
($\approx$\$421{,}900/year), concentrated in legacy infrastructure rather than the
live product. When the credits lapsed, roughly \textbf{\$27{,}783/month} of
previously credit-masked spend---about \$926/day---flipped from a hidden line to
cash. The footprint spanned \textbf{185 GCP projects} (27 billing-enabled, only
$\sim$4 actually needed) across \textbf{11 billing accounts}:

\begin{table}[h]
\centering
\small
\begin{tabular}{p{0.46\textwidth} r r}
\toprule
\textbf{Billing account} & \textbf{\$/month} & \textbf{\$/year} \\
\midrule
Legacy / startup-credit account (the waste) & 27{,}783 & 333{,}396 \\
Production account A & 4{,}040 & 48{,}480 \\
Secondary product group & 3{,}087 & 37{,}044 \\
New-stack non-prod (already efficient) & 248 & 2{,}976 \\
\midrule
\textbf{Total (actual billing)} & \textbf{$\sim$35{,}158} & \textbf{$\sim$421{,}900} \\
\bottomrule
\end{tabular}
\caption{Actual monthly GCP spend by billing account, taken from the billing
console after startup credits expired (11 billing accounts in total; the four
material ones shown). A single legacy account accounted for $\sim$79\% of the total;
the current-generation stack was already a rounding error.}
\end{table}

The legacy spend bought capacity the new design no longer needed: more than
\textbf{125 GKE nodes across 20 clusters} (single clusters of 19, 17, and 14 nodes;
64 of those nodes were 18\,GB machines), two standing \textbf{GPU virtual machines}
(an A100-class instance at $\sim$\$2{,}900/month and an L4-class instance at
$\sim$\$800/month, $\sim$\$3{,}700/month together), \textbf{5 Cloud SQL PostgreSQL
instances}, and $\sim$\textbf{18.5\,TB} of persistent disk---several terminated VMs
were even still paying for idle attached storage. Two problems compound here:
\textbf{cost} (most of the spend bought unused capacity) and \textbf{sovereignty}
(the entire footprint lived inside one hyperscaler---GKE, Cloud SQL, BigQuery, GCS,
Artifact Registry, Cloud Build---with the data gravity and lock-in that implies).

% ============================================================================
\section{The Migration: Two Documented Steps}
% ============================================================================

The migration consolidated the workload onto Hanzo's open-source stack---the same
self-hostable components (identity, key management, data, gateway, agent runtime)
that run sovereign and disconnected---and onto owned hardware for the
compute-heavy services. The savings landed in two measured steps:

\begin{table}[h]
\centering
\small
\begin{tabular}{p{0.50\textwidth} r r}
\toprule
\textbf{State} & \textbf{\$/month} & \textbf{vs.\ start} \\
\midrule
Actual billing (credits expired) & $\sim$35{,}158 & --- \\
\midrule
\emph{Step 1 --- retire legacy:} stop 2 GPU VMs ($-\$3{,}700$),
delete legacy GKE clusters, close dead projects & $\sim$6{,}489 & $-82\%$ \\
\emph{Step 2 --- consolidate \& right-size} onto the owned stack
and owned hardware & $\sim$3{,}500 & $-90\%$ \\
\midrule
\emph{Target --- finish consolidation} (retire last legacy projects) & $\sim$1{,}200 & $\approx$$-97$ pts \\
\bottomrule
\end{tabular}
\caption{Reduction to date. Step 1 (retiring unused legacy) alone removed $\sim$82\%
of spend with no impact on the live product; Step 2 consolidated the remaining
workload onto the self-hostable stack, landing at $\sim$\$3{,}500/month
($\approx$\$380{,}000/year saved). The $\sim$\$1{,}200/month line is a forward
\emph{target} as the last legacy projects are retired---not yet realized.}
\end{table}

Decommissioning was clean precisely because the legacy capacity was unused: shutting
down the two GPU VMs removed $\sim$\$3{,}700/month on its own; the largest legacy
GKE clusters and dozens of dead projects accounted for most of the rest. The net
effect is a reduction from $\sim$\$35{,}158 to $\sim$\$3{,}500 per month
($\approx$90\%) and, equally important, the elimination of single-vendor lock-in:
the resulting stack runs on owned hardware, in another cloud, or air-gapped.

% ============================================================================
\section{Where the Savings Come From}
% ============================================================================

The reduction is not a one-time clean-up; it is structural, and it reproduces:

\begin{itemize}
  \item \textbf{No metered model markup.} Owned models (the Zen family) replace
  per-token API calls, removing a recurring per-inference charge.
  \item \textbf{No hyperscaler margin.} Self-hosting on owned hardware removes the
  cloud provider's margin on compute, storage, and egress.
  \item \textbf{Efficient transport.} The ZAP protocol's 91\% memory and 96\%
  energy reductions and 30$\times$ throughput mean the same work needs far less
  hardware (see the ZAP paper).
  \item \textbf{Consolidation.} One coherent stack replaces a sprawl of projects,
  clusters, and services---the single largest line item above was pure redundant
  overhead.
\end{itemize}

Together these are the mechanism behind the roughly order-of-magnitude
infrastructure savings, and they apply to any organization carrying a metered-API
plus hyperscaler-sprawl cost structure.

% ============================================================================
\section{The Build: Consolidation, Not Volume}
% ============================================================================

The economics of \emph{running} the platform are only half the story; the economics
of \emph{building} it are the other half---and here the measured story is
\emph{doing more with dramatically less code}. The legacy estate was a sprawl of
\textbf{200+ microservices} (a single legacy auth monolith alone fronted 391 route
handlers; $\sim$80+ distinct deployable services, 49 of them active Node.js
processes at the time of migration). The same agentic engineering stack that now
runs the platform also built it, and the artifacts are small: \textbf{3 compiled Go
binaries} (ATS / broker-dealer / transfer-agent) alongside $\sim$10 OSS platform
services and 5 validator nodes. Measured git-over-git, the active codebase fell from
\textbf{3{,}071{,}930} to \textbf{1{,}056{,}566} lines ($\approx$70\% less; roughly
2.0\,M lines of dead code deleted):

\begin{table}[h]
\centering
\small
\begin{tabular}{p{0.40\textwidth} r r r}
\toprule
\textbf{Measure} & \textbf{Legacy} & \textbf{New stack} & \textbf{Change} \\
\midrule
Deployable services & 200+ & 3 Go binaries & $\sim$60$\times$ fewer \\
Active lines of code (whole estate) & 3{,}071{,}930 & 1{,}056{,}566 & $\approx$70\% less \\
Lines of code (trading/auth core) & $\sim$263{,}803 & $\sim$44{,}308 & $\sim$6$\times$ less \\
Vendored dependencies (flagship svc) & 1.3\,GB & 0 & eliminated \\
Provisioned cluster RAM & $\sim$1\,TB+ & $\sim$50\,MB working set & $\sim$99.5\% less \\
Deploy artifact (flagship) & --- & 48\,MB binary & replaces 17+ services \\
Deploy time & $\sim$1\,hr manual & $<$2\,min rolling & --- \\
Cold start (per service) & $\sim$10\,s & $<$1\,s & $\sim$10$\times$ faster \\
Automated tests (new stack) & n/a tracked & 700+ Go tests & --- \\
\bottomrule
\end{tabular}
\caption{Measured engineering footprint, new stack vs.\ the legacy implementation it
replaced, git-verified. The whole-estate line count fell from $\sim$3.07\,M to
$\sim$1.06\,M active lines; the trading/auth core specifically went $\sim$6$\times$
smaller. A single 48\,MB Go binary subsumes 17+ formerly separate services.}
\end{table}

This inverts the usual ``big system = big codebase'' intuition: agent-driven
engineering produced a \emph{leaner} system than the one it replaced, which is why
it is cheaper to run (fewer moving parts), cheaper to maintain (less code), and
faster to deploy ($<$2\,min rolling deploys, versus $\sim$1\,hr manual or 15--30\,min
canary before). The memory story is the same shape: the legacy fleet held
$\sim$1\,TB+ of \emph{provisioned} cluster RAM (64 nodes $\times$ 18\,GB) against an
application working set of roughly 10\,GB; the consolidated Go binary holds its
working set in $\sim$50\,MB ($\approx$99.5\% less), and cold start dropped from
$\sim$10\,s per Node.js service to $<$1\,s. On security, an internal source audit
found \textbf{93 findings} across 127+ repositories (22 critical, 30 high, 32
medium, 9 low) plus $\sim$780+ npm dependency vulnerabilities; the consolidated
rebuild remediated that entire class to \emph{zero actionable} at the architecture
level (header hardening, strict CORS allowlisting, JWKS-verified JWTs, secrets moved
out of plaintext environment, and the 1.3\,GB vendored dependency tree---the source
of most of the npm vulns---eliminated outright).

% ============================================================================
\section{The Compliance Bar}
% ============================================================================

The load-bearing proof point is that this lean, cheap-to-run system is not a
prototype: it is a regulated securities venue---an SEC-registered alternative
trading system with broker-dealer and transfer-agent functions---carried through
\textbf{FINRA/SEC approval} and \textbf{SOC\,2} compliance, in record time, with a
\textbf{small team}. Agent-driven engineering shipped and passed regulatory scrutiny
in financial services, among the most compliance-heavy software environments that
exists.

% ============================================================================
\section{Implications for Regulated and Government Modernization}
% ============================================================================

The same motion transfers directly to any regulated-enterprise or government
modernization---a healthcare system, a bank, a utility, or a government
contractor (defense included). An organization with an existing, audited (and,
for government systems, already-accredited) software stack does not need to
rebuild from zero; the agentic stack can modernize that stack in place---to a
post-quantum, sovereign, self-hostable footprint---while the consolidation cuts
recurring infrastructure cost by an order of magnitude, shrinks the codebase that
must be maintained (and, where applicable, accredited), and removes single-vendor
dependency. The wedge
is low-friction (modernize what you already have, with proofs), the savings are
immediate and structural, and the result is software the operator owns and can run
where the data lives. A regulated financial platform is the existence proof; the
same pattern applies anywhere the cost structure and the compliance burden repeat.

\paragraph{A note on figures.} Cost figures are drawn from the client's actual GCP
billing-console data after promotional credits expired (a more defensible basis than
list-price estimates); infrastructure, code, and security figures are drawn from
internal infrastructure, code, and source-security audits. Line counts, service
counts, and dependency sizes are measured (git-verified) where stated. Where a
number is a forward target rather than a realized result (e.g.\ the $\sim$\$1{,}200/
month consolidation floor) it is labeled as such; where a source labels a number
projected or self-reported, we have not relied on it here. All identifiers are
removed; the client is not named.

\end{document}
